% GENERATED FILE. Edit canonical lessons, then run npm run generate:books.
% canonical-source-hash: fnv1a64:c5306e0a3e8c2289
% canonical-lessons: SA-C285-adhinam, SA-C285-samayadhinam, SA-C285-yadi-samayah-abhavisyat, SA-W285-six-lines, SA-R285-depends-and-would

\chapter{It Depends, and What Would Have Been}
\label{ch:sa-it-depends-and-what-would-have-been}
\hlchaptermodality{285}

\begin{chapteropening}
\textbf{By the end of this chapter:} \emph{I can say what something depends on (tat samayādhīnam) and what would have happened (yadi samayaḥ abhaviṣyat, tarhi aham āgamiṣyam), and tell a short trip in six connected lines.}

Complete SA-R285-depends-and-would: say what something depends on and what would have happened.
\end{chapteropening}

\section[\texorpdfstring{\sk{अधीनम्} (adhīnam)}{अधीनम् (adhīnam)}]{\sk{अधीनम्} (adhīnam) — depending on}

\label{lesson:SA-C285-adhinam}



\begin{quote}
Before the new one: say \textquotedblleft{}if there is time, I will come\textquotedblright{}. (\emph{yadi samayaḥ bhaviṣyati, tarhi aham āgamiṣyāmi}.) A shorter answer says only what it depends on.
\end{quote}

\subsection*{You'll want to know: \sk{अधीनम्}}
\textbf{\sk{अधीनम्}} — \emph{adhīnam} — \textquotedblleft{}depending on\textquotedblright{}. It comes at the end of what it depends on, joined to it as one word: \emph{samaya} and \emph{adhīnam} meet as \textbf{\sk{समयाधीनम्}} (\emph{samayādhīnam}).

\begin{quote}
\textbf{\sk{तत्} \sk{समयाधीनम्।}}
\end{quote}

\begin{quote}
\emph{tat samayādhīnam.} — \textquotedblleft{}That depends on time.\textquotedblright{}
\end{quote}

\begin{quote}
\textbf{\sk{सर्वं} \sk{वर्षाधीनम्।}}
\end{quote}

\begin{quote}
\emph{sarvaṃ varṣādhīnam.} — \textquotedblleft{}Everything depends on the rain.\textquotedblright{}
\end{quote}

\begin{quote}
\textbf{\sk{एतत्} \sk{धनाधीनम्।}}
\end{quote}

\begin{quote}
\emph{etat dhanādhīnam.} — \textquotedblleft{}This depends on money.\textquotedblright{}
\end{quote}

\begin{cousinweb}
\textbf{\sk{अधीनम्}} (\emph{adhīnam}) is built on \textbf{\sk{अधि}} (\emph{adhi}), \textquotedblleft{}over, on\textquotedblright{}: what rests on something. Monier-Williams gives \emph{adhīna} as \textquotedblleft{}resting on\textquotedblright{}, and \textquotedblleft{}depending on, subject to\textquotedblright{}, always at the end of a compound, and Macdonell gives it the same way, which is why this book joins it to the word before it. Pali kept it as \emph{adhīna}, \textquotedblleft{}dependent on\textquotedblright{} (CDIAL 266).
\end{cousinweb}

\subsection*{Your turn}
Say these aloud:

\begin{itemize}
\raggedright
  \item \emph{adhīnam}
  \item \textquotedblleft{}That depends on time\textquotedblright{} — \emph{tat samayādhīnam}
  \item \textquotedblleft{}That depends on the rain\textquotedblright{} — \emph{tat varṣādhīnam}
\end{itemize}

\begin{tcolorbox}[breakable,colback=teal!4,colframe=teal!35!black,title={Before you move on}]
What does \textbf{\sk{अधीनम्}} mean? (\textquotedblleft{}Depending on\textquotedblright{}.) \textquotedblleft{}It depends on money\textquotedblright{}? (\textbf{\sk{तत्} \sk{धनाधीनम्}}.)
\end{tcolorbox}
\section[\texorpdfstring{\sk{तत्} \sk{समयाधीनम्} (tat samayādhīnam)}{तत् समयाधीनम् (tat samayādhīnam)}]{\sk{तत्} \sk{समयाधीनम्} (tat samayādhīnam) — it depends}

\label{lesson:SA-C285-samayadhinam}



\begin{quote}
Say \textquotedblleft{}that depends on time\textquotedblright{}. (\emph{tat samayādhīnam}.) Then say \textquotedblleft{}if it rains, I stay at home\textquotedblright{}. (\emph{yadi varṣā bhavati, tarhi ahaṃ gṛhe tiṣṭhāmi}.)
\end{quote}

\subsection*{You'll want to know: what depends sets the ending}
\textbf{-\sk{अधीन}} describes the thing that depends, so it takes that thing's gender: \emph{-aḥ} for a masculine noun, \emph{-ā} for a feminine one, \emph{-am} for a neuter one.

\begin{quote}
\textbf{\sk{मम} \sk{निर्णयः} \sk{धनाधीनः।}}
\end{quote}

\begin{quote}
\emph{mama nirṇayaḥ dhanādhīnaḥ.} — \textquotedblleft{}My decision depends on money.\textquotedblright{}
\end{quote}

\begin{quote}
\textbf{\sk{यात्रा} \sk{वर्षाधीना।}}
\end{quote}

\begin{quote}
\emph{yātrā varṣādhīnā.} — \textquotedblleft{}The trip depends on the rain.\textquotedblright{}
\end{quote}

\begin{quote}
\textbf{\sk{कार्यं} \sk{समयाधीनम्।}}
\end{quote}

\begin{quote}
\emph{kāryaṃ samayādhīnam.} — \textquotedblleft{}The work depends on time.\textquotedblright{}
\end{quote}

\textbf{\sk{निर्णयः}} is masculine, so \textbf{-\sk{अधीनः}}; \textbf{\sk{यात्रा}} is feminine, so \textbf{-\sk{अधीना}}.

\subsection*{You'll want to know: it depends, and then the condition}
Say what it depends on, then lay out the condition with \textbf{\sk{यदि}} … \textbf{\sk{तर्हि}}:

\begin{quote}
\textbf{\sk{मम} \sk{यात्रा} \sk{वर्षाधीना।}}
\end{quote}

\begin{quote}
\emph{mama yātrā varṣādhīnā.} — \textquotedblleft{}My trip depends on the rain.\textquotedblright{}
\end{quote}

\begin{quote}
\textbf{\sk{यदि} \sk{वर्षा} \sk{न} \sk{भविष्यति},}
\end{quote}

\begin{quote}
\emph{yadi varṣā na bhaviṣyati,}
\end{quote}

\begin{quote}
\textbf{\sk{तर्हि} \sk{वयं} \sk{गमिष्यामः।}}
\end{quote}

\begin{quote}
\emph{tarhi vayaṃ gamiṣyāmaḥ.} — \textquotedblleft{}If it does not rain, we will go.\textquotedblright{}
\end{quote}

\begin{quote}
\textbf{\sk{तत्} \sk{भवदधीनम्।}}
\end{quote}

\begin{quote}
\emph{tat bhavadadhīnam.} — \textquotedblleft{}That is up to you.\textquotedblright{}
\end{quote}

\textbf{\sk{भवदधीनम्}} (\emph{bhavadadhīnam}) is \emph{bhavat}, the stem of \textbf{\sk{भवान्}}, with \emph{adhīnam}: \textquotedblleft{}depending on you\textquotedblright{}, that is, \textquotedblleft{}up to you\textquotedblright{}.

\subsection*{Your turn}
Say these aloud:

\begin{itemize}
\raggedright
  \item \textquotedblleft{}My decision depends on money\textquotedblright{} — \emph{mama nirṇayaḥ dhanādhīnaḥ}
  \item \textquotedblleft{}My decision depends on time\textquotedblright{} — \emph{mama nirṇayaḥ samayādhīnaḥ}
  \item \textquotedblleft{}My decision depends on the rain\textquotedblright{} — \emph{mama nirṇayaḥ varṣādhīnaḥ}
\end{itemize}

\begin{tcolorbox}[breakable,colback=teal!4,colframe=teal!35!black,title={Before you move on}]
\textquotedblleft{}The trip depends on the rain\textquotedblright{}? (\textbf{\sk{यात्रा} \sk{वर्षाधीना}}.) Why \textbf{-\sk{अधीना}}? (\textbf{\sk{यात्रा}} is feminine.) \textquotedblleft{}That is up to you\textquotedblright{}? (\textbf{\sk{तत्} \sk{भवदधीनम्}}.)
\end{tcolorbox}
\section[\texorpdfstring{\sk{यदि} \sk{समयः} \sk{अभविष्यत्} (yadi samayaḥ …)}{यदि समयः अभविष्यत् (yadi samayaḥ …)}]{\sk{यदि} \sk{समयः} \sk{अभविष्यत्} (yadi samayaḥ abhaviṣyat) — what would have been}

\label{lesson:SA-C285-yadi-samayah-abhavisyat}



\begin{quote}
Say \textquotedblleft{}if there is time, I will come\textquotedblright{}. (\emph{yadi samayaḥ bhaviṣyati, tarhi aham āgamiṣyāmi}.) Then say \textquotedblleft{}he went\textquotedblright{}, the old way. (\emph{saḥ agacchat}.) The next form puts those two together.
\end{quote}

\subsection*{You'll want to know: the conditional}
\begin{quote}
\textbf{\sk{यदि} \sk{समयः} \sk{अभविष्यत्},}
\end{quote}

\begin{quote}
\emph{yadi samayaḥ abhaviṣyat,}
\end{quote}

\begin{quote}
\textbf{\sk{तर्हि} \sk{अहम्} \sk{अगमिष्यम्।}}
\end{quote}

\begin{quote}
\emph{tarhi aham agamiṣyam.} — \textquotedblleft{}If there had been time, I would have gone.\textquotedblright{}
\end{quote}

\textbf{\sk{अभविष्यत्}} (\emph{abhaviṣyat}), \textquotedblleft{}would have been\textquotedblright{}, is the future stem \textbf{\sk{भविष्य}-} with the \textbf{\sk{अ}-} and the \emph{-t} of \textbf{\sk{अगच्छत्}} around it. For \textquotedblleft{}I\textquotedblright{} the ending is \emph{-am}: \textbf{\sk{अगमिष्यम्}} (\emph{agamiṣyam}), \textquotedblleft{}I would have gone\textquotedblright{}. It says that the thing did not happen (Pāṇini 3.3.139; the \textbf{\sk{अ}-}, 6.4.71).

\noindent
\begin{tabularx}{\linewidth}{@{}>{\raggedright\arraybackslash}X>{\raggedright\arraybackslash}X@{}}
\toprule
\textbf{the future} & \textbf{the conditional} \\
\midrule
\textbf{\sk{भविष्यति}}, will be & \textbf{\sk{अभविष्यत्}}, would have been \\
\textbf{\sk{गमिष्यति}}, will go & \textbf{\sk{अगमिष्यत्}}, would have gone \\
\textbf{\sk{करिष्यति}}, will do & \textbf{\sk{अकरिष्यत्}}, would have done \\
\textbf{\sk{गमिष्यामि}}, I will go & \textbf{\sk{अगमिष्यम्}}, I would have gone \\
\bottomrule
\end{tabularx}

After the prefix \textbf{\sk{आ}-}, the \textbf{\sk{अ}-} melts into its long \emph{ā}: \textbf{\sk{आगमिष्यम्}} (\emph{āgamiṣyam}), \textquotedblleft{}I would have come\textquotedblright{}.

\begin{quote}
\textbf{\sk{यदि} \sk{समयः} \sk{अभविष्यत्},}
\end{quote}

\begin{quote}
\emph{yadi samayaḥ abhaviṣyat,}
\end{quote}

\begin{quote}
\textbf{\sk{तर्हि} \sk{अहम्} \sk{आगमिष्यम्।}}
\end{quote}

\begin{quote}
\emph{tarhi aham āgamiṣyam.} — \textquotedblleft{}If there had been time, I would have come.\textquotedblright{}
\end{quote}

\subsection*{You'll want to know: both halves, and \textquotedblleft{}if he were to\textquotedblright{}}
Both halves take it, for a past that did not come about:

\begin{quote}
\textbf{\sk{यदि} \sk{धनम्} \sk{अभविष्यत्},}
\end{quote}

\begin{quote}
\emph{yadi dhanam abhaviṣyat,}
\end{quote}

\begin{quote}
\textbf{\sk{तर्हि} \sk{सः} \sk{कार्यम्} \sk{अकरिष्यत्।}}
\end{quote}

\begin{quote}
\emph{tarhi saḥ kāryam akariṣyat.} — \textquotedblleft{}If there had been money, he would have done the work.\textquotedblright{}
\end{quote}

You may also hear \textbf{\sk{यदि}} with the \emph{-et} of advice, \textquotedblleft{}if he were to …\textquotedblright{}; meet it, and say the conditional yourself:

\begin{quote}
\textbf{\sk{यदि} \sk{सः} \sk{आगच्छेत्},}
\end{quote}

\begin{quote}
\emph{yadi saḥ āgacchet,}
\end{quote}

\begin{quote}
\textbf{\sk{तर्हि} \sk{मीरा} \sk{प्रसन्ना} \sk{भवेत्।}}
\end{quote}

\begin{quote}
\emph{tarhi mīrā prasannā bhavet.} — \textquotedblleft{}If he were to come, Mira would be pleased.\textquotedblright{}
\end{quote}

\subsection*{Your turn}
Say these aloud:

\begin{itemize}
\raggedright
  \item \textquotedblleft{}If there had been time, I would have gone\textquotedblright{} — \emph{yadi samayaḥ abhaviṣyat, tarhi aham agamiṣyam}
  \item \textquotedblleft{}If there had been time, he would have gone\textquotedblright{} — \emph{yadi samayaḥ abhaviṣyat, tarhi saḥ agamiṣyat}
  \item \textquotedblleft{}If there had been time, he would have done it\textquotedblright{} — \emph{yadi samayaḥ abhaviṣyat, tarhi saḥ akariṣyat}
\end{itemize}

\begin{tcolorbox}[breakable,colback=teal!4,colframe=teal!35!black,title={Before you move on}]
\textquotedblleft{}Would have been\textquotedblright{}? (\textbf{\sk{अभविष्यत्}}.) \textquotedblleft{}I would have come\textquotedblright{}? (\textbf{\sk{आगमिष्यम्}}.) Did the thing happen? (No: that is what the conditional says.)
\end{tcolorbox}
\section[(six lines)]{Six lines about a trip}

\label{lesson:SA-W285-six-lines}



\begin{quote}
This task asks for six lines that tell a story: a trip, real or invented.

Pick one reader: Mira, or a friend you name. Say \textbf{\sk{घटितम्}}, \textbf{\sk{दृढः}}, \textbf{\sk{औषधम्}} and \textbf{\sk{ईदृशः}} aloud; one of them will go in.
\end{quote}

\subsection*{What to know first}
A told story is not a list of events. It has \textbf{events} in the past of \emph{-vān} and \emph{-vatī}, kept in order with \textbf{\sk{प्रथमम्}}, \textbf{\sk{ततः}}, \textbf{\sk{अनन्तरम्}} and \textbf{\sk{अन्ते}}, and joined with \emph{-tvā} or \emph{-ya}. It has a \textbf{scene}, with \textbf{\sk{स्म}}, or with \textbf{\sk{यदा}} … \textbf{\sk{तदा}}. It has a \textbf{reason}, with \textbf{\sk{यतः}} or \textbf{-\sk{कारणात्}}, and a \textbf{feeling}, with \textbf{\sk{अहं} \sk{प्रसन्नः} \sk{अस्मि}} or \textbf{\sk{प्रसन्ना}}.

A good one ends by looking at what might have been: \textbf{\sk{यदि}} … \textbf{\sk{अभविष्यत्}}, \textbf{\sk{तर्हि}} ….

\subsection*{Writing — connected composition}
On paper, in Devanagari, write \textbf{six lines} to your reader about the trip, one line for each of these:

\begin{enumerate}
\raggedright
  \item what happened, and when: \textbf{\sk{ह्यः} \sk{अहं} … \sk{गतवान्।}} or \textbf{\sk{गतवती।}};
  \item the scene as it began: \textbf{\sk{यदा} \sk{अहं} … \sk{स्म}, \sk{तदा} …};
  \item what happened next, in order: \textbf{\sk{प्रथमम्}} or \textbf{\sk{ततः}}, with a \emph{-tvā} or \emph{-ya} form, and one thing that went wrong: \textbf{… \sk{नष्टः}} or \textbf{\sk{धूमयानं} \sk{गतम्}};
  \item why, with \textbf{\sk{यतः} …} or \textbf{…\sk{कारणात्}};
  \item how it ended and how you felt, with \textbf{\sk{अन्ते}} and \textbf{\sk{अहम्}} … \textbf{\sk{अस्मि}}: \textbf{\sk{प्रसन्नः}}, \textbf{\sk{श्रान्ता}}, \textbf{\sk{उत्सुकः}};
  \item what would have been, with \textbf{\sk{यदि} … \sk{अभविष्यत्}, \sk{तर्हि} …}, and one word with \textbf{\sk{घ}}, \textbf{\sk{ढ}}, \textbf{\sk{औ}} or \textbf{\sk{ई}}: \textbf{\sk{घटितम्}}, \textbf{\sk{दृढः}}, \textbf{\sk{औषधम्}}, \textbf{\sk{ईदृशः}}.
\end{enumerate}

There is no model in this lesson. Use only words and forms this book has taught, written whole, with the one join you know: \textbf{\sk{गृहं} \sk{गच्छामि}}. The work is choosing the right past for each line, and joining the lines.

\begin{tcolorbox}[breakable,colback=teal!4,colframe=teal!35!black,title={Before you move on}]
Now read your lines as you read Mira's day. Mark every verb \textbf{e} for an event and \textbf{s} for the scene: line two \textbf{s}, lines one and three mostly \textbf{e}.

Then three checks. A past in \emph{-vān} or \emph{-vatī} follows the one who did it: \textbf{\sk{अहं} \sk{गतवती}}, if you are a woman. A verb with a prefix takes \emph{-ya}: \textbf{\sk{प्रविश्य}}, \textbf{\sk{उपविश्य}}. A feeling takes your gender: \textbf{\sk{प्रसन्नः}} or \textbf{\sk{प्रसन्ना}}.

Last, read line six alone. Does its \textbf{\sk{यदि}} half say what did not happen, with \textbf{\sk{अभविष्यत्}}?
\end{tcolorbox}
\section[(review)]{Review — it depends, and what would have been}

\label{lesson:SA-R285-depends-and-would}



\begin{quote}
Say \textquotedblleft{}if he does not come, I go\textquotedblright{}, and \textquotedblleft{}if you are ill, please go to the doctor\textquotedblright{}. (\emph{yadi saḥ na āgacchati, tarhi ahaṃ gacchāmi}; \emph{yadi bhavān rugṇaḥ asti, tarhi vaidyaṃ gacchatu}.) Then \textquotedblleft{}otherwise\textquotedblright{}. (\emph{anyathā}.)
\end{quote}

\subsection*{Your turn}
Say these aloud:

\begin{itemize}
\raggedright
  \item That depends on time.
  \item The trip depends on the rain.
  \item My decision depends on money.
  \item If there had been time, I would have come.
  \item If there had been money, he would have done the work.
\end{itemize}

\begin{tcolorbox}[breakable,colback=teal!4,colframe=teal!35!black,title={Before you move on}]
\textbf{\sk{वर्षाधीना}}: what depends, a masculine or a feminine thing? (A feminine one, such as \textbf{\sk{यात्रा}}.) \textbf{\sk{अभविष्यत्}}: did it happen? (No.) Now the other way round: \textquotedblleft{}I would have gone\textquotedblright{}? (\textbf{\sk{अगमिष्यम्}}.) \textquotedblleft{}That is up to you\textquotedblright{}? (\textbf{\sk{तत्} \sk{भवदधीनम्}}.)
\end{tcolorbox}
